\documentclass[11pt]{article}

%%% Packages
\usepackage{amsmath,amsthm,amssymb,amsfonts}
\usepackage{mathtools}
\usepackage{enumerate}
\usepackage{hyperref}
\usepackage{url}
\usepackage{booktabs}
\usepackage{array}
\usepackage{xcolor}
\usepackage{graphicx}
\usepackage[margin=1in]{geometry}

%%% Theorem environments
\newtheorem{theorem}{Theorem}[section]
\newtheorem{lemma}[theorem]{Lemma}
\newtheorem{proposition}[theorem]{Proposition}
\newtheorem{corollary}[theorem]{Corollary}
\newtheorem{conjecture}[theorem]{Conjecture}
\theoremstyle{definition}
\newtheorem{definition}[theorem]{Definition}
\newtheorem{example}[theorem]{Example}
\newtheorem{remark}[theorem]{Remark}

%%% Notation
\newcommand{\Z}{\mathbb{Z}}
\newcommand{\Q}{\mathbb{Q}}
\newcommand{\R}{\mathbb{R}}
\newcommand{\C}{\mathbb{C}}
\newcommand{\N}{\mathbb{N}}
\newcommand{\HH}{\mathbb{H}}
\DeclareMathOperator{\tr}{tr}
\DeclareMathOperator{\SL}{SL}
\DeclareMathOperator{\AUC}{AUC}

\title{Continued Fraction Invariants of Three-Body Periodic Orbits:\\
The Golden Ratio Structure of the Figure-Eight\\
and a Periodic Table of 695 Orbits}

\author{
  Saar Shai\thanks{Independent Researcher. Email: \texttt{saar.shai@gmail.com}.}
  \and
  Claude Opus 4.6\thanks{AI research assistant by Anthropic; contributed to mathematical analysis, exact-arithmetic computation, and manuscript preparation.}
}

\date{March 2026}

\begin{document}

\maketitle

\begin{abstract}
Planar three-body periodic orbits are classified topologically by
braid words in the free group~$F_2$. Using the classical isomorphism
$F_2 \cong \Gamma(2) \subset \SL(2,\Z)$, each orbit determines a
$2 \times 2$ integer matrix whose attracting fixed point is a
quadratic irrational with an eventually periodic continued fraction
(CF). We apply exact quadratic-surd arithmetic to all 695
equal-mass orbits in the Li--Liao catalog, correcting the 17.7\%
CF corruption rate of naive floating-point computation. For the
figure-eight orbit, the fixed-point equation is exactly
$z^2 - z - 1 = 0$, the minimal polynomial of the golden
ratio~$\varphi$; this identity is forced by the Fibonacci structure
of the commutator matrix in~$\Gamma(2)$ and is
convention-independent (trace~$= 18$, discriminant~$= 320 = 64
\times 5$). Across the full catalog, CF ``nobility'' (fraction of
partial quotients equal to~$1$) anticorrelates with braid entropy
at partial~$\rho = -0.890$ ($p < 10^{-50}$, controlling for word
length) and predicts low-entropy orbits with $\AUC = 0.980$ on a
held-out test set of 195 orbits, providing a microsecond-cost
screening tool for orbit discovery.  \textbf{Caveat:} braid entropy
measures topological complexity of the symbolic coding, not physical
dynamical stability; confirming the latter requires Floquet
multiplier computations not yet available for this catalog. We organize all 691 orbits with detected CF
periods into a $9 \times 8$ ``periodic table'' indexed by CF
period length and geometric mean, revealing structure orthogonal to
topological family classification (69\% purity per cell) and
identifying empty cells as predictions for undiscovered orbit
types. While the $F_2 \to \Gamma(2) \to \text{CF}$ bridge is
classical and the KAM-theoretic connection between noble
frequencies and stability is well known, our contribution is the
systematic exact-arithmetic quantification across the largest
available catalog, extending Kin--Nakamura--Ogawa's
Stern--Brocot framework from ${\sim}\!13$ Lissajous orbits to
695 orbits across all families.
\end{abstract}

\medskip
\noindent\textbf{Keywords:} three-body problem, periodic orbits, continued fractions, Farey tessellation, golden ratio, braid groups, Stern--Brocot tree, KAM theory.

\noindent\textbf{MSC 2020:} 70F07, 37J40, 11A55, 37E30, 70H12.

%% ================================================================
\section{Introduction}\label{sec:intro}
%% ================================================================

\subsection{Periodic orbits of the three-body problem}

The gravitational three-body problem remains one of the central
unsolved problems in classical mechanics. While the general problem
is chaotic and admits no closed-form solution, \emph{periodic}
orbits play a distinguished role: by the Birkhoff--Lewis theorem
they are dense in the phase space, and they organize the global
dynamics as the skeleton around which quasi-periodic and chaotic
trajectories are arranged.

The modern era of three-body periodic orbit discovery began with
Moore's~\cite{Moore1993} numerical discovery of the figure-eight
orbit and its existence proof by Chenciner and
Montgomery~\cite{ChencinerMontgomery2000}. Large-scale numerical
searches by \v{S}uvakov and Dmitra\v{s}inovi\'{c}~\cite{Suvakov2013}
and Li and Liao~\cite{LiLiao2017, LiLiao2019} have produced catalogs
of hundreds to thousands of periodic orbits for equal-mass systems
with zero angular momentum.

\subsection{Topological classification via braids}

Montgomery~\cite{Montgomery1998} observed that three-body orbits in
the plane can be classified by their \emph{braid type}: the
worldlines of three bodies in spacetime $(x, y, t)$ form a
three-strand braid. For planar orbits avoiding triple collision, the
relevant structure is the fundamental group of the shape sphere
(configuration space modulo translations and rotations, with collision
points removed), which is isomorphic to the free group $F_2$ on two
generators.

The Li--Liao catalog assigns each orbit a free-group word in the
generators $\{a, b, A, B\}$ (where $A = a^{-1}$, $B = b^{-1}$),
encoding how the orbit's projection to the shape sphere winds
around the two-body collision singularities.

\subsection{Why number theory?}

The free group $F_2$ is isomorphic to the principal congruence
subgroup $\Gamma(2) \subset \SL(2,\Z)$ via the standard generators
\begin{equation}\label{eq:generators}
a \mapsto \begin{pmatrix} 1 & 2 \\ 0 & 1 \end{pmatrix}, \qquad
b \mapsto \begin{pmatrix} 1 & 0 \\ 2 & 1 \end{pmatrix}.
\end{equation}
Each word $w \in F_2$ thus determines a $2 \times 2$ integer matrix
$M(w) \in \Gamma(2)$ acting on the upper half-plane~$\HH$ by
M\"obius transformations. When $|\tr(M)| > 2$ (the hyperbolic case,
which holds for all 695 orbits in the catalog), $M$ has two real
fixed points, both quadratic irrationals. Their continued fraction
expansions are eventually periodic (by Lagrange's theorem) and
encode arithmetic information about the orbit.

This perspective connects three-body dynamics to the \emph{Farey
tessellation} of~$\HH$: the tessellation by ideal triangles whose
vertices are Farey fractions, invariant under $\SL(2,\Z)$. The
cutting sequence of the geodesic connecting the two fixed points of
$M(w)$ through the Farey tessellation directly yields the CF
expansion. This framework was developed by
Series~\cite{Series1985} and applied to three-body braids by
Kin, Nakamura, and Ogawa~\cite{KNO2023}, who classified Lissajous
three-braids via the Stern--Brocot tree for approximately 13 orbits.

\subsection{Our contribution}

We extend the Kin--Nakamura--Ogawa framework from ${\sim}\!13$
Lissajous orbits to the full Li--Liao catalog of 695 equal-mass,
zero-angular-momentum periodic orbits. Our specific contributions
are:

\begin{enumerate}
\item \textbf{An exact anchor result:} the figure-eight orbit maps
  to the golden ratio~$\varphi$ under~$\Gamma(2)$, an algebraic
  identity forced by the Fibonacci structure of the commutator
  matrix (Section~\ref{sec:figure-eight}).

\item \textbf{Exact arithmetic at scale:} we replace floating-point
  CF computation (which corrupts 17.7\% of orbits) with exact
  quadratic-surd arithmetic, obtaining correct periodic CFs for 691
  of 695 orbits (Section~\ref{sec:catalog}).

\item \textbf{A CF periodic table:} we organize all 691 orbits into
  a $9 \times 8$ grid indexed by CF period length and geometric
  mean, revealing structure orthogonal to topological classification
  (Section~\ref{sec:periodic-table}).

\item \textbf{Quantified predictive power:} CF nobility achieves
  $\AUC = 0.980$ for blind prediction of low-entropy orbits,
  providing a microsecond-cost screening tool
  (Section~\ref{sec:prediction}).

\item \textbf{Honest accounting of limitations:} we carefully
  distinguish algebraic correlations (CF vs.\ trace of the
  \emph{same} matrix) from physical ones (CF vs.\ orbital period
  $T^*$ from numerical integration), and state what would be needed
  for full physical validation (Section~\ref{sec:discussion}).
\end{enumerate}

We emphasize that the mathematical framework ($F_2 \cong \Gamma(2)$,
CF of quadratic irrationals, Farey tessellation) is entirely
classical. The novelty lies in the systematic, exact-arithmetic
application across the largest available catalog.


%% ================================================================
\section{Mathematical Framework}\label{sec:framework}
%% ================================================================

\subsection{From free-group words to $\Gamma(2)$ matrices}

Let $F_2 = \langle a, b \rangle$ be the free group on two
generators. The map $\phi: F_2 \to \Gamma(2)$ defined
by~\eqref{eq:generators} is an isomorphism onto
$\Gamma(2) = \ker(\SL(2,\Z) \to \SL(2,\Z/2\Z))$.
The inverses are
\[
\phi(A) = \begin{pmatrix} 1 & -2 \\ 0 & 1 \end{pmatrix}, \qquad
\phi(B) = \begin{pmatrix} 1 & 0 \\ -2 & 1 \end{pmatrix}.
\]
For a word $w = g_1 g_2 \cdots g_n$ with $g_i \in \{a, b, A, B\}$,
the matrix is $M(w) = \phi(g_1)\phi(g_2) \cdots \phi(g_n)$.

\subsection{Fixed points and quadratic irrationals}

The matrix $M = \bigl(\begin{smallmatrix} \alpha & \beta \\
\gamma & \delta \end{smallmatrix}\bigr)$ acts on
$\hat{\R} = \R \cup \{\infty\}$ by the M\"obius transformation
$T(z) = (\alpha z + \beta)/(\gamma z + \delta)$. When
$|\tr(M)| = |\alpha + \delta| > 2$, the transformation is
\emph{hyperbolic} with two real fixed points:
\begin{equation}\label{eq:fixed-points}
z_{\pm} = \frac{(\alpha - \delta) \pm
\sqrt{(\alpha + \delta)^2 - 4}}{2\gamma}.
\end{equation}
The discriminant $\Delta = \tr(M)^2 - 4$ is a positive non-square
integer for all 695 orbits, so the fixed points are quadratic
irrationals in $\Q(\sqrt{\Delta})$.

The \emph{attracting} fixed point~$z_+$ (where $|T'(z_+)| < 1$) is
the one we associate to the orbit. By Lagrange's theorem, every
quadratic irrational has an eventually periodic CF expansion.

\begin{remark}[Invariance]
The trace $\tr(M)$ is a conjugacy-class invariant: cyclic
permutations of the word yield conjugate matrices with the same
trace and discriminant.
\end{remark}

\subsection{Continued fraction invariants}

Every quadratic irrational~$\alpha$ has a CF expansion
$\alpha = [a_0; a_1, a_2, \ldots]$ that is eventually periodic.
We define the following CF invariants for each orbit:
\begin{itemize}
\item \textbf{CF period length}~$\ell$: the minimal period of the
  eventually periodic CF.
\item \textbf{CF geometric mean}
  $G = \bigl(\prod_{i=k}^{k+\ell-1} a_i\bigr)^{1/\ell}$:
  geometric mean of partial quotients in one period.
\item \textbf{CF nobility}
  $\nu = \#\{i \in [k, k+\ell): a_i = 1\} / \ell$:
  fraction of partial quotients equal to~$1$.
\end{itemize}

We use ``nobility'' by analogy with noble numbers in KAM theory.
\textbf{This is a non-standard extension} of the classical binary
(noble/not-noble) concept; the standard measures in KAM theory are
the Diophantine exponent and Brjuno sum.

\subsection{The Farey tessellation and cutting sequences}

The \emph{Farey tessellation} of the upper half-plane~$\HH$ is the
ideal triangulation whose edges connect Farey neighbors $p/q$ and
$r/s$ (with $|ps - qr| = 1$) by hyperbolic geodesics.  This
tessellation is invariant under $\SL(2,\Z)$, and the quotient
$\HH/\Gamma(2)$ is the thrice-punctured sphere---which is also the
shape sphere of the planar three-body problem with collision points
removed~\cite{Montgomery1998}.

For a hyperbolic element $M \in \Gamma(2)$ with attracting and
repelling fixed points $z_+, z_-$, the \emph{cutting sequence} of the
geodesic from $z_-$ to $z_+$ through the Farey tessellation encodes
the CF expansion of~$z_+$: each time the geodesic crosses $a_i$
consecutive edges of one type, the CF records the partial
quotient~$a_i$.  This gives a geometric realization of the CF: the
Farey tessellation is the geometric substrate, and the CF is the
combinatorial record of how the orbit's axis traverses
it~\cite{Series1985}.

\subsection{Braid entropy}

For a hyperbolic matrix~$M$ with $|\tr(M)| > 2$, the \emph{braid
entropy} is
\begin{equation}\label{eq:entropy}
h(w) = \frac{\log \lambda(M)}{|w|}
\end{equation}
where $\lambda(M)$ is the spectral radius and $|w|$ is the word
length.

\textbf{Important caveat:} This is the entropy of the symbolic
coding (the braid), not the Lyapunov exponent of the physical
orbit. The correlation between braid entropy and true dynamical
stability, while expected on general grounds, remains to be verified
by direct computation of Floquet multipliers.


%% ================================================================
\section{The Figure-Eight and the Golden Ratio}\label{sec:figure-eight}
%% ================================================================

\subsection{The computation}

The figure-eight orbit has the free-group word $BabA$, which is a
cyclic rotation of the commutator $[a,b] = abAB$. Computing the
$\Gamma(2)$ matrix:
\begin{equation}\label{eq:fig-eight-matrix}
M(BabA) = \phi(B)\phi(a)\phi(b)\phi(A) =
\begin{pmatrix} 5 & -8 \\ -8 & 13 \end{pmatrix}.
\end{equation}
This matrix has trace~$18$, determinant~$1$, and discriminant
$\Delta = 18^2 - 4 = 320 = 64 \times 5$. The fixed points satisfy
$\frac{5z - 8}{-8z + 13} = z$, which simplifies to
\begin{equation}\label{eq:golden-poly}
z^2 - z - 1 = 0,
\end{equation}
the \emph{minimal polynomial of the golden ratio}
$\varphi = (1+\sqrt{5})/2$.

\subsection{Why the golden ratio is structurally forced}

The appearance of~$\varphi$ is not a numerical coincidence. The
matrix entries $\{5, 8, 13\}$ are consecutive Fibonacci numbers
$F_5, F_6, F_7$, a consequence of the commutator structure. The
characteristic polynomial reduces to $z^2 - z - 1 = 0$ precisely
because the Fibonacci relation is encoded in the matrix. The
factor of~$5$ in the discriminant $320 = 64 \times 5$ guarantees
that $\sqrt{5}$ --- and hence $\varphi$ --- appears in every
fixed-point computation, regardless of generator conventions. All
fixed points across all conjugate representatives are golden-ratio
algebraic conjugates living in $\Q(\sqrt{5})$.

\subsection{Verification}

The figure-eight fixed point was verified to 60 decimal places
using arbitrary-precision arithmetic. The CF expansion
$\varphi = [1; \overline{1}]$ was verified to 30 terms. This is an
exact algebraic identity.


%% ================================================================
\section{Continued Fraction Properties of the Full Catalog}\label{sec:catalog}
%% ================================================================

\subsection{Data and methodology}

We analyzed all 695 equal-mass, zero-angular-momentum periodic
orbits from the Li--Liao catalog~\cite{LiLiao2017, LiLiao2019}.

\textbf{High-precision arithmetic.} Naive floating-point computation of CFs
degrades catastrophically for orbits with long words: 123 of 695
orbits (17.7\%) had corrupted CF data with float64 arithmetic. We
replaced this with exact integer quadratic-surd tracking: since every
fixed point has the form $(P + q\sqrt{D})/Q$ with integers
$P, q, Q, D$, the CF can be computed by tracking the integer state
$(P_n, q_n, Q_n)$ at each step using exact integer arithmetic, with
\texttt{mpmath} at 100 digits used only for floor evaluation. This
yielded correct periodic CFs for 691 of 695 orbits (99.4\%); the
remaining 4 have CF periods exceeding 300 terms.

The improvement was dramatic: the partial correlation
$\rho(\text{nobility}, \text{braid entropy})$ strengthened from
$-0.538$ (float64) to $-0.890$ (exact).

\subsection{Catalog-wide correlations}

All 695 orbits are hyperbolic ($|\tr(M)| > 2$). The following are
Spearman rank correlations; partial correlations control for word
length via OLS residualization.

\begin{table}[ht]
\centering
\caption{Key correlations ($N = 695$, exact surd arithmetic).}
\label{tab:correlations}
\begin{tabular}{llrl}
\toprule
Relationship & Type & $\rho$ & Note \\
\midrule
Nobility vs.\ braid entropy & Partial (ctrl.\ $|w|$)
  & $-0.890$ & (a) \\
Nobility vs.\ CF period length & Partial (ctrl.\ $|w|$)
  & $-0.962$ & (a) \\
Nobility vs.\ $T^*$ & Raw & $-0.535$ & (b) \\
Nobility vs.\ $\log|\tr|$ & Raw & $-0.546$ & (a) \\
Word length vs.\ $T^*$ & Raw & $+0.999$ & Baseline \\
\bottomrule
\end{tabular}

\medskip
\small
Notes: (a)~Braid entropy and $\log|\tr|$ are computed from the
\emph{same} $\Gamma(2)$ matrix as CF properties; these correlations
are partly algebraic. (b)~$T^*$ is the physical period from numerical
orbit integration, independent of the matrix.
\end{table}

All reported correlations survive permutation testing (10,000
permutations, $p < 0.0001$). Within-family correlations remain
strong ($\rho \sim -0.85$ for families I.A, I.B, II.C
individually).

\subsection{Family structure}

The Kruskal--Wallis test detects significant differences in mean
nobility across topological families ($H = 16.36$, $p = 0.0026$).

\begin{table}[ht]
\centering
\caption{CF properties by topological family.}
\label{tab:families}
\begin{tabular}{lrrrr}
\toprule
Family & $N$ & Mean nobility & Mean CF geom.\ mean & CF period mode \\
\midrule
I.A   & 190 & 0.713 & 1.358 & 14 \\
I.B   & 191 & 0.686 & 1.391 & 5 \\
II.A  & 4   & 0.921 & 1.066 & 1 \\
II.B  & 10  & 0.755 & 1.368 & 1 \\
II.C  & 300 & 0.732 & 1.406 & varies \\
\bottomrule
\end{tabular}
\end{table}


%% ================================================================
\section{A Periodic Table of Three-Body Orbits}\label{sec:periodic-table}
%% ================================================================

The CF invariants naturally suggest a two-dimensional organization,
analogous to the chemical periodic table. We construct a
$9 \times 8$ grid in which \textbf{rows} correspond to CF period
length (binned) and \textbf{columns} correspond to the geometric
mean of the periodic CF partial quotients (binned from $1.00$
to~$1.45$). Rows index algebraic complexity; columns index nobility
(a braid-entropy proxy, not a direct physical stability measure).

Of the 72 cells, 51 are populated and 21 are empty. The
figure-eight orbit occupies the (period~$1$, gmean $1.00$--$1.05$)
cell \emph{alone} among all 691 orbits --- it is ``hydrogen'' in
the analogy. Its CF is $[0; \overline{1}]$, the simplest possible.

\textbf{Orthogonality to topological classification.} The average
family purity per cell is only 0.688, meaning the dominant family
in a typical cell accounts for just 69\% of its orbits. The CF
organization captures structure \textbf{orthogonal} to the
topological family classification.

\textbf{Empty-cell predictions.} Among the 21 empty cells, several
are surrounded by populated neighbors, suggesting orbits with those
CF properties may exist but have not yet been cataloged.  These
constitute tentative predictions: orbits with these specific
combinations of algebraic complexity and nobility could be sought
in future numerical searches.

\textbf{Noble gases.}  For each row (CF period class), the orbit
with the highest nobility serves as the analog of a noble gas---the
lowest braid-entropy element of its period.  These orbits systematically have
$G \approx 1.0$, meaning their CF partial quotients are predominantly
$1$s (Fibonacci-like structure).  As the CF period increases, the
noble-gas nobility converges toward the golden-ratio limit, connecting
braid entropy to Fibonacci structure quantitatively.


%% ================================================================
\section{Predictive Power}\label{sec:prediction}
%% ================================================================

\subsection{Blind classification of low-entropy orbits}

We split the 695 orbits into training (500) and test (195) sets.
The target is whether braid entropy falls in the bottom 30th
percentile.

\begin{table}[ht]
\centering
\caption{Blind prediction results (test set, $N = 195$).}
\label{tab:prediction}
\begin{tabular}{lrr}
\toprule
Model features & Accuracy & AUC \\
\midrule
Word length only & 0.677 & 0.741 \\
Nobility only & 0.974 & 0.980 \\
Word length + nobility & 0.959 & 0.985 \\
Full CF features & 0.979 & 0.991 \\
\bottomrule
\end{tabular}
\end{table}

\textbf{Critical caveat.} The prediction target (braid entropy) is
computed from the same $\Gamma(2)$ matrix as the CF features. This
test demonstrates that CF properties efficiently encode information
about matrix trace growth --- algebraically interesting but not the
same as predicting physical stability. A fully physical validation
would require Floquet multipliers from linearized orbit integration.

\subsection{Nobility-guided orbit search}

As a more physically grounded test, we compared two strategies for
finding near-periodic solutions (return errors below a threshold of
$1.0$--$2.0$ in position):
\begin{itemize}
\item \textbf{Random search:} uniform sampling of initial conditions
  $(v_1, v_2) \in [-1,1]^2$.
\item \textbf{Nobility-guided search:} generate $\Gamma(2)$ elements
  of word length $2$--$8$, rank by CF nobility, and focus numerical
  integration around the most noble candidates.
\end{itemize}
In a preliminary experiment (200 trials), the nobility-guided search
found approximately twice as many near-periodic orbits as random
search at loose return-error thresholds ($1.0$--$2.0$).  At the
stricter threshold of~$0.5$, the advantage narrowed to a factor
of~${\sim}\!1.06$, suggesting nobility provides coarse guidance
toward the right region of phase space but does not substitute for
fine numerical refinement.

\subsection{Nobility as a microsecond screening tool}

Computing CF nobility requires only integer matrix multiplication
and exact CF expansion via the quadratic-surd algorithm (microseconds).
In contrast, computing braid entropy from a physical trajectory
requires hours of numerical integration. At $\AUC = 0.980$, the
nobility filter correctly ranks 98\% of orbit pairs by relative
braid entropy, eliminating the need to numerically integrate
${\sim}\!70\%$ of high-entropy candidates.

We propose the following screening pipeline for large-scale orbit
discovery: (1)~generate candidate braid words of target complexity;
(2)~compute nobility via exact surd CF (cost: microseconds per
candidate); (3)~rank candidates and apply a nobility threshold;
(4)~invest in full numerical integration only for candidates
passing the threshold.

The screening tool's effectiveness depends critically on exact surd
computation: with float64 arithmetic, $\AUC$ drops to $0.853$.


%% ================================================================
\section{Discussion}\label{sec:discussion}
%% ================================================================

\subsection{What this means}

The central finding is that the classical
$F_2 \to \Gamma(2) \to \text{CF}$ correspondence, applied
systematically, reveals a strong and quantifiable relationship
between CF structure and orbital complexity. The figure-eight
anchors the picture at one extreme: it maps to the golden ratio
with the simplest CF. More complex orbits map to less noble fixed
points.

\subsection{Connection to KAM theory}

The qualitative picture is consistent with KAM theory
(Kolmogorov~\cite{Kolmogorov1954}, Arnold~\cite{Arnold1963},
Moser~\cite{Moser1962}): orbits whose frequency ratios are
well-approximated by rationals (resonant, low nobility) are less
stable, while those with noble frequency ratios resist
resonance-driven instability.  Our quantitative finding---that
nobility explains significant variance in braid entropy even after
controlling for word length---suggests that the CF structure encodes
information beyond mere orbit complexity.

However, \emph{braid entropy is not the same as dynamical
instability}.  The braid entropy measures topological complexity of
the symbolic coding; the dynamical Lyapunov exponent measures
divergence of nearby physical trajectories.  Confirming the physical
KAM connection requires Floquet-multiplier computations that the
current catalog does not provide.

\subsection{Limitations}

\begin{enumerate}
\item \textbf{Algebraic vs.\ physical correlations.} The braid
  entropy is computed from the same matrix as the CF properties.
  Strong correlations are therefore partly algebraic.

\item \textbf{No physical stability data.} The catalog provides
  orbital periods but not Floquet multipliers or Lyapunov exponents.

\item \textbf{Gap predictions are not validated.} Direct $N$-body
  integration of top predictions found 0 strictly periodic orbits.

\item \textbf{Non-standard ``nobility'' measure.} Our continuous
  $\nu \in [0,1]$ is highly correlated ($|\rho| = 0.994$) with the
  CF geometric mean.

\item \textbf{Catalog dependence.} All results depend on the
  correctness of the Li--Liao catalog's word assignments.
\end{enumerate}

\subsection{Future directions}

\begin{enumerate}
\item Compute Floquet multipliers for the Li--Liao catalog and test
  whether CF nobility predicts linear stability.
\item Extend to unequal-mass three-body problems.
\item Correlate Hamilton's action with CF properties.
\item Target the empty cells in the periodic table with focused
  numerical searches.
\item Investigate analogous structure for the $N$-body problem
  ($N > 3$) in higher braid groups.
\end{enumerate}


%% ================================================================
\section{Conclusion}\label{sec:conclusion}
%% ================================================================

The classical isomorphism $F_2 \cong \Gamma(2)$ provides a concrete
bridge between three-body periodic orbits and number theory.  The
figure-eight orbit maps exactly to the golden ratio~$\varphi =
(1+\sqrt{5})/2$, with the fixed-point equation $z^2 - z - 1 = 0$
forced by the Fibonacci structure of the commutator matrix
$\bigl(\begin{smallmatrix}5&-8\\-8&13\end{smallmatrix}\bigr)$---not
a numerical coincidence.  Across all 695 orbits in the Li--Liao
catalog (691 with exact CF periods), continued-fraction properties
computed with exact quadratic-surd arithmetic correlate strongly with
braid entropy (partial $\rho = -0.890$, $\AUC = 0.980$ on a held-out
test set of 195 orbits from a 500/195 train/test split), survive
permutation testing, and organize
the catalog into a $9 \times 8$ periodic table (51 populated cells)
that reveals structure orthogonal to topological family
classification.

The golden ratio's appearance is the algebraic shadow of the
simplest commutator in~$\Gamma(2)$, projected through the Farey
tessellation.  Its quantitative reach across 695 orbits---predicting
braid entropy (a topological, not physical, measure) at $\AUC = 0.98$
from continued fractions alone on a held-out test set,
organizing the full catalog into a periodic table that reveals
structure invisible to topological classification, and identifying
empty cells as candidates for undiscovered orbits---suggests that
the number theory of quadratic irrationals has more to say about
celestial mechanics than has yet been fully explored.  Whether these
algebraic correlations extend to physical dynamical stability remains
an open question requiring Floquet multiplier data.


%% ================================================================
\section*{Acknowledgments}
%% ================================================================

We thank Xiaoming Li and Shijun Liao for making their three-body
periodic orbit catalog publicly available, and Eiko Kin, Hidetoshi
Nakamura, and Mitsuhiko Ogawa for the foundational work connecting
three-body braids to the Stern--Brocot tree.


\section*{Data and Code Availability}

The Li--Liao three-body periodic orbit catalog is available at
\url{https://github.com/sjtu-liao/three-body}. Our computational
pipeline and all 695 orbit CF decompositions will be made available
at \url{https://github.com/SaarShai/Primes-Equispaced} upon
publication.


%% ================================================================
\begin{thebibliography}{99}

\bibitem{Arnold1963}
V.~I.~Arnold,
``Small denominators and problems of stability of motion in
classical and celestial mechanics,''
\emph{Russian Math.\ Surveys}, \textbf{18}(6):85--191, 1963.

\bibitem{ChencinerMontgomery2000}
A.~Chenciner and R.~Montgomery,
``A remarkable periodic solution of the three-body problem in the
case of equal masses,''
\emph{Ann.\ Math.}, \textbf{152}(3):881--901, 2000.

\bibitem{KNO2023}
E.~Kin, H.~Nakamura, and M.~Ogawa,
``Lissajous 3-braids,''
\emph{J.\ Math.\ Soc.\ Japan}, \textbf{75}(1):195--228, 2023.

\bibitem{Kolmogorov1954}
A.~N.~Kolmogorov,
``On the conservation of conditionally periodic motions under small
perturbation of the Hamiltonian,''
\emph{Dokl.\ Akad.\ Nauk SSSR}, \textbf{98}:527--530, 1954.

\bibitem{LiLiao2017}
X.~Li and S.~Liao,
``More than six hundred new families of Newtonian periodic planar
collisionless three-body orbits,''
\emph{Sci.\ China Phys.\ Mech.\ Astron.},
\textbf{60}(12):129511, 2017.

\bibitem{LiLiao2019}
X.~Li and S.~Liao,
``Collisionless periodic orbits in the free-fall three-body
problem,''
\emph{New Astronomy}, \textbf{70}:22--26, 2019.

\bibitem{Montgomery1998}
R.~Montgomery,
``The $N$-body problem, the braid group, and action-minimizing
periodic solutions,''
\emph{Nonlinearity}, \textbf{11}(2):363, 1998.

\bibitem{Moore1993}
C.~Moore,
``Braids in classical dynamics,''
\emph{Phys.\ Rev.\ Lett.}, \textbf{70}(24):3675--3679, 1993.

\bibitem{Moser1962}
J.~Moser,
``On invariant curves of area-preserving mappings of an annulus,''
\emph{Nachr.\ Akad.\ Wiss.\ G\"ottingen}, II:1--20, 1962.

\bibitem{Series1985}
C.~Series,
``The modular surface and continued fractions,''
\emph{J.\ London Math.\ Soc.}, \textbf{31}(1):69--80, 1985.

\bibitem{Suvakov2013}
M.~\v{S}uvakov and V.~Dmitra\v{s}inovi\'{c},
``Three classes of Newtonian three-body planar periodic orbits,''
\emph{Phys.\ Rev.\ Lett.}, \textbf{110}(11):114301, 2013.

\end{thebibliography}

\appendix

\section{Exact Quadratic-Surd CF Algorithm}\label{app:surd}

For a fixed point $z = (P + q\sqrt{D})/Q$ with integers
$P, q, Q, D$ (where $D$ is not a perfect square), the CF is
computed using exact integer quadratic surd representation, with
\texttt{mpmath} at 100 digits used only for floor evaluation:
\begin{enumerate}
\item Compute $a_n = \lfloor z_n \rfloor$ using \texttt{mpmath}
  at 100-digit precision for the floor of the surd.
\item Update $z_{n+1} = 1/(z_n - a_n)$, yielding a new quadratic
  surd with integer coefficients via denominator rationalization
  (this step is exact integer arithmetic).
\item Record the state $(P_n, q_n, Q_n)$; the first repetition
  detects the exact period.
\end{enumerate}
The surd tracking is exact; only the floor computation uses
floating-point (at 100 digits, far beyond the precision needed).
This eliminates the CF corruption that affected 17.7\% of orbits
in naive float64 implementations.

\end{document}
